\documentclass[11pt]{article}
\usepackage[margin=1in]{geometry}
\usepackage{amsmath,amssymb,amsthm,mathtools}
\usepackage[hidelinks]{hyperref}
\allowdisplaybreaks[2]
\setlength{\emergencystretch}{3em}
\newtheorem{proposition}{Proposition}[section]
\newtheorem{lemma}[proposition]{Lemma}
\theoremstyle{remark}\newtheorem{remark}[proposition]{Remark}
\newcommand{\E}{\mathbb E}
\newcommand{\dd}{\,d}
\newcommand{\one}{\mathbf1}
\newcommand{\Gp}{G_+}
\title{Global matrix-value envelope, signed gate excess,\\
and the two-sided Q3 split}
\author{Analytical increment for review}
\date{Version 1; October 22, 2026}
\begin{document}\maketitle
\begin{abstract}
The proposed matrix-value comparison closes the off-chart
angular calculation without a receiver derivative. Its constant
is explicit in the integrated weak-diagonal excess. A signed
identity bounds that excess by the positive output moment,+$\frac12\int(Qc+N_-)$, without a pointwise smallness assumption
on $N_-$. The stated stopped guards do not supply this integrated
budget: CLOCK's full passage forcing budget is distinct from
the c-form remainder budget, and CFR's negative-part estimate
concerns a different moment. Thus no uniform global-envelope
or full S2 theorem is claimed. Fixed residual-guard dependence
of an eventual distance constant is admissible in the tail
consumers. Separately, Q3 labels away from two initial strips
of width $O(q\sqrt{\log(1/q)})$ remain nearly frozen through
the stopped endpoint. Their integrated mass divided by $q$
is $O(q^9\log(1/q))$. The two strips' entry obligations are
stated without assuming recruitment or changing prior cohorts.
\end{abstract}

\section{Setting and scope of the four tasks}
All new labels have prefix \texttt{an03gesgv1:}.
We use DECG's original balanced flow, not a comparison trajectory:
\[
 P=\tfrac12P_1+\tfrac12P_2,\quad
 P_1=N(e_1,q^2I_2),\quad P_2=N(\rho e_2,q^2I_2),\quad
 \rho\in[13/20,7/10],\quad \lambda=\rho^2,\quad Q=\lambda+q^2,
\]
\[
 R_\theta=\E_P[(X-f)X^T\one_{A_\theta}],\quad
 A_\theta=\{a_\theta\cdot X>0\},\quad
 U=\sqrt m\,a_\theta,\quad W=\sqrt m\,a_\psi,
 \quad a_v=(\cos v,\sin v).
\]
The exact equations are
\begin{equation}
 \theta'=a_\psi^TR_\theta a_\theta^\perp,\qquad
 \psi'=(a_\psi^\perp)^TR_\theta a_\theta,\qquad
 (\log m)'=2a_\psi^TR_\theta a_\theta.
 \label{an03gesgv1:exact}
\end{equation}
Write $e_j=X_j-f_j$, $c=Q^{-1}\E_{P_2}[f_2X_2]$,
$I(t)=Q\int_{t_0}^t(1-c)$, $\beta=\lambda/Q$, and
$\ell=\log(1/q)$. Work on the stated stopped interval,
$t_0=t_h$, $t_*-t_0\le L_T\ell$, with
\begin{equation}
 0\le c\le c_b<1,\quad \int_{t_0}^{t_*}r<B,\quad
 r=\sup_\theta|R_{11}|,\quad
 \|e_2\|_{P_1}+\|e_1\|_{P_2}\le M_{\rm res}q,
 \quad Y<K_Y,\quad E_{\Gp}<E_g.
 \label{an03gesgv1:guards}
\end{equation}
Loss monotonicity gives a fixed bound $C_{\mathcal L}$ on
\emph{all} component residual $L^2(P_i)$ norms, also before
$t_0$. The cross-residual guard improves two of these bounds
on the stopped interval. Hence, as in BULK and DECG,
\begin{equation}
 |R_{12}|+|R_{21}|\le a_*q,\qquad
 a_*:=C_R(1+M_{\rm res}),
 \label{an03gesgv1:cross}
\end{equation}
where $C_R$ depends only on the loss bound and ratio box.

The same-field reference is never reset. Assume its stopped
bound $|\theta_c|,|\psi_c|\le qR_c$, $R_c=79/51$, and
the fixed-chart c-form hypothesis used by DECG. To distinguish
two genuinely different quantities, write
\[
 F_F:=\int_{t_0}^{t_*}\epsilon_F\le\overline F_F,
 \qquad F_{\rm pass}:=\int_{t_0}^{t_*}f_*.
\]
Only the first budget is part of the inherited DECG field
guard. The second is an assumption in CLOCK's passage theorem;
reading that theorem does not supply the assumption.
Section~\ref{an03gesgv1:budgetsection} also records what follows
if a later bootstrap explicitly guards $F_{\rm pass}$.

Put $D=(|\theta-\theta_c|+|\psi-\psi_c|)/q$, $\xi=D(t_0)$.
Angles use continuous lifts on the strong branch. Before the
first $qD=\delta_0$, take $0<\delta_0<1/4$ and $qR_c\le1/4$;
then $|\theta|,|\psi|\le1/2$. These comparisons concern strong
complement and SEP labels initially on that branch, not Q3
angles near $\pi$. Their fixed S1 entry caps give
$q\xi<\delta_0$ for small $q$. Task 4 uses a separate argument for Q3.

\section{Task 1: the global matrix-value comparison}
Define the nonnegative, receiver-uniform excess
\begin{equation}
 d_2=I'/2,\qquad
 X(t)=\sup_\theta\bigl(|R_{22}(t,\theta)|-d_2(t)\bigr)_+,
 \quad \mathfrak X(s,t)=\int_s^t X.
 \label{an03gesgv1:X}
\end{equation}
It is finite on compact time intervals by the residual bounds.
No uniform-in-$q$ bound on its time integral is presumed.

\begin{proposition}[Global sine envelope and Duhamel estimate]
\label{an03gesgv1:sine}
For every receiver, $N=|\sin\theta|+|\sin\psi|$ satisfies
\begin{equation}
 D^+N\le (r+d_2+X+a_*q)N+a_*q.
 \label{an03gesgv1:N}
\end{equation}
Set $a_0=\lambda(1-c_b)/2>0$ and
$\varepsilon_q=(a_*q+q^2/2)L_T\ell=o(1)$ for fixed guards.
For any $t_0\le s\le t\le t_*$,
\begin{equation}
 \frac{N(t)}q\le
 e^{\beta[I(t)-I(s)]/2+\int_s^tr+\mathfrak X(s,t)+\varepsilon_q}
 \left[\frac{N(s)}q+\frac{a_*}{a_0}\right].
 \label{an03gesgv1:duhamel}
\end{equation}
While both endpoint angles are on the physical strong branch
specified above,
\begin{equation}
 D(t)\le C_\delta
 e^{\beta[I(t)-I(s)]/2+\int_s^tr+\mathfrak X(s,t)+\varepsilon_q}
 [D(s)+W_*],\quad
 C_\delta=\sec(1/2),\quad W_*=4R_c+a_*/a_0.
 \label{an03gesgv1:globalD}
\end{equation}
\end{proposition}
\begin{proof}
Expanding \eqref{an03gesgv1:exact}, and writing
$v=|\sin\theta|$, $w=|\sin\psi|$, gives
\begin{align*}
 D^+v&\le |R_{11}|v+|R_{22}|w+|R_{12}|+|R_{21}|vw,\\
 D^+w&\le |R_{11}|w+|R_{22}|v+|R_{21}|+|R_{12}|vw.
\end{align*}
These upper Dini bounds also hold at a zero of a sine.
Since $vw\le v+w$, summation proves \eqref{an03gesgv1:N}.
No receiver derivative of $R_\theta$ occurs.

Keep the $Q$--$\lambda$ difference explicitly:
$d_2=\beta I'/2+q^2(1-c)/2$.
Variation of constants for $N/q$, with the nonnegative
coefficient $r+X+a_*q+q^2(1-c)/2$ separated out, bounds
the inhomogeneous integral by
\[
 a_*\int_s^t e^{-\beta[I(v)-I(s)]/2}\dd v
 \le a_*\int_s^t e^{-a_0(v-s)}\dd v\le a_*/a_0.
\]
The remaining small coefficient integrates to at most
$\varepsilon_q$. This proves \eqref{an03gesgv1:duhamel}.

For $|\theta|,|\psi|\le1/2$, the mean value theorem gives
$|\theta|\le C_\delta|\sin\theta|$ and the analogous
bound for $\psi$. Also sine is globally 1-Lipschitz. Thus
\[
 D\le C_\delta N/q+2R_c,\qquad N/q\le D+2R_c.
\]
Insert these inequalities at the two endpoints of
\eqref{an03gesgv1:duhamel}, and use that its exponential
is at least one, to obtain \eqref{an03gesgv1:globalD}.
Comparability would be false for arbitrary angle lifts near
$\pi$; that case is not used in this argument.
\end{proof}

\begin{proposition}[One chart exit suffices; reentry causes no loss]
\label{an03gesgv1:splice}
Choose a fixed $L_c>R_c$ and put $d_L=L_c-R_c$.
Let
\[
 C_{\rm loc}=p_*^{-1}e^{C_T\overline F_F},\quad
 p_* =\Phi(-7L_c/10),\quad
 C_T=2+p_*^{-1}+(7/10)\phi(0)/p_*,
\]
as in DECG Proposition 2.1. For every originally charted or
initially exterior strong-branch receiver, up to physical exit,
\begin{equation}
 D(t)\le \mathcal C_D(t_*)\,\xi e^{\beta I(t)/2},\qquad
 \mathcal C_D(t_*)=
 C_\delta C_{\rm loc}
 e^{B+\mathfrak X(t_0,t_*)+\varepsilon_q}
       (1+W_*/d_L).
 \label{an03gesgv1:spliced}
\end{equation}
This is an estimate with its unpaid time integral visible,
not a fixed-constant global-envelope theorem.
\end{proposition}
\begin{proof}
Until the first exit $\tau_L$ from the scaled square use
DECG's proved local lemma. If there is no exit, it already
implies \eqref{an03gesgv1:spliced}. At an exit,
$D(\tau_L)\ge L_c-R_c=d_L$. Consequently
$D(\tau_L)+W_*\le(1+W_*/d_L)D(\tau_L)$.
Apply \eqref{an03gesgv1:globalD} from that time onwards and
substitute the local bound at $\tau_L$. This estimate holds
through every later reentry; it is not restarted at each chart
crossing. An initially exterior receiver has $\xi\ge d_L$,
so apply the same argument from $t_0$, without a local stage.
Coincidence with the reference is covered by the local lemma.
The equality case $L_c=R_c$ supplies no separation $d_L$;
one can take, for example, $L_c=R_c+1$ instead. No derivative
contract beyond this fixed square is used.
\end{proof}

\section{Task 2: signed gate excess and its unpaid budget}
\begin{lemma}[The signed cancellation is valid without the proposed cap]
\label{an03gesgv1:signed}
Put $N_-=\E_{P_2}[(f_2X_2)_-]$, with $z_-=(-z)_+$.
For every receiver, exactly
\begin{equation}
 R_{22}-d_2=\tfrac12\E_{P_1}[e_2X_2\one_A]
              -\tfrac12\E_{P_2}[e_2X_2\one_{A^c}].
 \label{an03gesgv1:identity}
\end{equation}
The following inequalities hold:
\begin{align}
 R_{22}-d_2
 &\le \tfrac12 M_{\rm res}q^2+
             \tfrac12\E_{P_2}[(f_2X_2)_+\one_{A^c}],
 \label{an03gesgv1:upper}\\
 -R_{22}-d_2
 &\le \tfrac12 M_{\rm res}q^2-	frac12Q(1-2c)+\tfrac12N_-.
 \label{an03gesgv1:lower}
\end{align}
In particular, without requiring $N_-\le Q(1-2c)$,
\begin{equation}
 X(t)\le\tfrac12M_{\rm res}q^2+\tfrac12(Qc+N_-),\qquad
 \mathfrak X(t_0,t_*)\le
 \tfrac Q2\int c+\tfrac12\int N_-+	frac12M_{\rm res}q^2L_T\ell.
 \label{an03gesgv1:Xbound}
\end{equation}
\end{lemma}
\begin{proof}
The global identity $\E_{P_2}[e_2X_2]=Q(1-c)$ proves
\eqref{an03gesgv1:identity}. The $P_1$ term has absolute
value at most $M_{\rm res}q^2/2$.
Since $e_2X_2=X_2^2-f_2X_2$, dropping the nonpositive
$-X_2^2$ term on $A^c$ proves \eqref{an03gesgv1:upper}.
Also
\[
 -R_{22}-d_2\le\tfrac12M_{\rm res}q^2
       +\tfrac12\E_{P_2}[(f_2X_2)_+]-d_2.
\]
The positive moment equals $Qc+N_-$, yielding
\eqref{an03gesgv1:lower}. Its right side is smaller than
$M_{\rm res}q^2/2+(Qc+N_-)/2$ by $Q(1-c)/2$.
Thus that latter nonnegative expression bounds both signs
of $R_{22}$ above $d_2$. Take the supremum and integrate.
The proposed pointwise restriction on $N_-$ is unnecessary.
\end{proof}

\subsection*{Which integral is actually available?}
\label{an03gesgv1:budgetsection}
CLOCK defines
$f_*=c+|V|+K_w+T+\epsilon_F$, with nonnegative terms in its
passage contract. Therefore an \emph{explicitly assumed}
$F_{\rm pass}\le\overline F_{\rm pass}$ would give
$\int c\le\overline F_{\rm pass}$ immediately, without
applying its passage theorem. The inherited remainder guard
$F_F\le\overline F_F$ does not give that inequality.
From the response guard alone the bound is only
\begin{equation}
 \int_{t_0}^{t_*}c\le c_bL_T\ell,
 \label{an03gesgv1:responseonly}
\end{equation}
which is not a fixed budget. No conclusion of CLOCK's
passage theorem is imported to supply its own input.

CFR's negative-part lemma controls
\[
 (K_O)_-,\qquad K_O=q^{-1}\E_{P_1}[(f_O)_2X_1],
 \qquad V_O=-(qQ)^{-1}\E_{P_2}[(f_O)_1X_2].
\]
Neither quantity is $N_-$: the former uses a signed integral
on $P_1$, the latter uses output component 1, while $N_-$
is the negative part \emph{before} integration of component
2 on $P_2$. Replacing $K_O$ by $(K_O)_+$ does not control it.
Nor does the signed moment $\E_{P_2}[f_2X_2]=Qc$ bound its
negative part. Loss gives only a pointwise $N_-\le C$, whose
direct time integral is $O(\ell)$.

One can expose the remaining cohort charges without a sign
guess. Whenever the instantaneous LCRC core geometry and
mass bound hold, RCET's positive $W_2$ on $\Gp$ gives
\begin{equation}
 N_-\le C\left[
 q^2\mu_{\mathcal K}+E_\Sigma+q\sqrt{E_\Sigma J_\Sigma}
       +\mu_{\mathcal B}
       +\mu_{\Gp}e^{-k_1/q^2}\right].
 \label{an03gesgv1:Ncohorts}
\end{equation}
Indeed $(\sum z_i)_-\le\sum(z_i)_-$, and
$f_{\Gp,2}\ge0$, so its negative product with $X_2$ is
supported on $X_2<0$. Gaussian tails pay that term.
For a general cohort,
$|f_{S,2}|\le2\sqrt{E_SJ_S}|X_1|+2E_S|X_2|$;
integration against $|X_2|$ proves the secondary charge.
Charted core donors give $Cq^2\mu_{\mathcal K}$ on $P_2$,
and $|f_{\mathcal B}|\le\mu_{\mathcal B}|X|$ pays the last
population. These arguments are state estimates only.
Under $J_\Sigma\le q^2K_Y^2$, the secondary part is
$E_\Sigma+q^2K_Y\sqrt{E_\Sigma}$, still unpaid by the
current guards. A bounded core mass and LCRC's retained
$\sup\mu_{\mathcal B}/q=o(1)$ would pay their contributions
to \emph{this} integral, since $q\ell=o(1)$; they do not
pay $\int\mu_{\mathcal B}/q$ in the field reduction.

The exact sufficient budget missing from this route is
\begin{equation}
 \boxed{\quad
 \int_{t_0}^{t_*}(Qc+N_-)\dd t\le B_{\rm gate}<\infty
 \quad\hbox{uniformly in }q,t_* .\quad}
 \label{an03gesgv1:missing}
\end{equation}
If the full CLOCK forcing integral is separately guarded,
the remaining item is precisely
\begin{equation}
 \boxed{\qquad\int_{t_0}^{t_*}N_-\dd t\le B_-<\infty
       \quad\hbox{uniformly in }q,t_* .\qquad}
 \label{an03gesgv1:missingnegative}
\end{equation}
A bounded constant, not $o(1)$, suffices for the distance
consumer. Neither boxed display is added to a claimed full
theorem. Task 2 stops at this budget check; accordingly
Task 1 has not produced a $q$-uniform $C_D$.
This is a missing estimate from the stated inputs, not an
assertion that an initialized trajectory violates it.

\section{Task 3: the constant audit}
An eventual finite budget in \eqref{an03gesgv1:missing} would
replace $\mathfrak X$ in \eqref{an03gesgv1:spliced} by a
fixed number. Denote the resulting distance constant by $C_D$.
The following audit concerns its downstream use; it does not
prove that budget.

\paragraph{Gain and secondary energy.}
DECG Corollary 3.3 uses
\[
 A_0=C_D^2\max(1,\delta_0^{-2}),\quad
 C_m=e^{6B+1/\beta+o(1)},\quad
 m(t,\alpha)\le C_m m(t_0,\alpha)
       \max(1,A_0q^2\xi_\alpha^2e^{\beta I})^{1+\nu},
 \quad \nu=q^2/\lambda.
\]
BULK Proposition 5.2 chooses a positive fixed threshold
$\eta_*<(\delta_0/(2C_D))^2$. Its energy constants depend
on $C_D,A_0,C_m,\eta_*^{-1}$ and the fixed entry cap. After
the SEP moments and stopped clock are inserted, these multiply
$q^5$ and $q^{\Gamma_\nu}H$, with
$\Gamma_\nu\ge289169/924000>3/10$.
There is no additional order-one weak energy term proportional
to $C_D$.

\paragraph{Strong-complement mass.}
In TAILREP's tail-moment lemma put
$A=A_0q^2e^{\beta I}$ in its scaled-displacement convention.
The result is proportional to
\[
 C_m\epsilon^p d_c^{-p}\max(1,A d_c^2)^{1+\nu},\qquad p=3/s.
\]
The fixed factor $A_0$ changes its constant, not its power of
$q$. The previously proved margin is
$\Delta_\nu\ge17141311/140024000>3/25$.
Consequently its prospective integrated field contribution is
$C(C_D,A_0,C_m)q^{\Delta_\nu}\ell=o(1)$.
No second displacement-growth factor is charged after this
mass calculation.

\paragraph{Return to LCRC.}
Its primary constant choices use the fixed core radius, core
mass bound, $B_s,c_b$ and model constants. The secondary and
complement terms just listed enter with positive powers of
$q$ (also after square roots and the fixed $K_Y$ factors in
CFR/LCRC). Thus a fixed $M_{\rm res}$ dependence of $C_D$
is admissible. For example the explicit
$1+W_*/(L_c-R_c)$ in \eqref{an03gesgv1:spliced} need not
be bounded independently of $M_{\rm res}$ before the guards
are chosen. Once they are fixed, decrease $q$ afterwards.
One must not replace $\mathfrak X$ by an unproved
$q$-dependent budget and then call its exponential a fixed
constant. That, rather than the fixed $M_{\rm res}$ factor,
is the current obstruction.

\paragraph{Chart choice and forcing smallness.}
No enlargement after the guards is needed for this audit;
$L_c=R_c+1$ suffices, and initially exterior labels are covered
by Proposition~\ref{an03gesgv1:splice}.
Alternatively one may choose a larger fixed $L_c$ after
$(B,E_g,K_Y,M_{\rm res})$ and before $C_0$, as permitted by
roadmap ledger 17. The c-form guard must hold on the chosen
square; validity on a smaller square does not itself extend it.
The constants $p_*^{-1},C_T$ then grow
but remain fixed. They enter the envelope only through
$C_{\rm loc}=p_*^{-1}e^{C_T\overline F_F}$.
Boundedness, not smallness of $C_T\overline F_F$, is enough
for these tail consumers. If a multiplier at most $e$ is
desired, the extra choice is $\overline F_F\le1/C_T$;
it is not required by the distance estimate.
CLOCK's separate passage condition is
$\overline F_{\rm pass}\le e_*/(4C_{\rm pass})$.
If a later construction uses a common budget for both
purposes, it may choose the minimum of these two bounds.
Any chart-dependent coefficient multiplying $C_0$ in CFR
is handled by choosing $C_0$ later. None of these choices
proves the unpaid forcing integrals. No order-one LCRC core
constant is redefined here.

\section{Task 4: lower-side no-gain and a frozen Q3 interior}
\begin{lemma}[Lower-side no-gain]
\label{an03gesgv1:lowernogain}
On any subinterval $[s,t]$ of the stopped interval on which
$\sin\theta<0$, independently of the sign of $\sin\psi$,
\begin{equation}
 2R_{22}\sin\theta\sin\psi\le Cq+CM_{\rm res}q^2,
 \qquad
 m(t)\le m(s)\exp\!\left(2\int_s^tr+
                   C(1+M_{\rm res})q(t-s)\right).
 \label{an03gesgv1:nogain}
\end{equation}
In particular the multiplier is at most $e^{2B+o(1)}$.
\end{lemma}
\begin{proof}
DECG's negative-gate estimate is
$|R_{22}|\le CM_{\rm res}q^2+Ce^{-k\sin^2\theta/q^2}$.
Multiply it by $2|\sin\theta\sin\psi|\le2|\sin\theta|$;
$\sup_{v\ge0}ve^{-kv^2/q^2}=q/\sqrt{2ke}$ proves the first
inequality. In \eqref{an03gesgv1:exact} the strong diagonal
costs $2r$ and cross terms cost $2a_*q$. Integrate.
The estimate applies while the input sine stays negative;
no assertion is made across a subsequent upper-side interval.
\end{proof}

\begin{lemma}[Both closed input gates: a stronger speed estimate]
\label{an03gesgv1:interiorspeed}
At any time, including before $t_0$, if
$\cos\theta\le-\Lambda q$ and $\sin\theta\le-\Lambda q$,
then
\begin{equation}
 \max_{i,j}|R_{ij}(\theta)|\le C_L e^{-k\Lambda^2},\qquad
 |\theta'|+|\psi'|+| (\log m)'|\le C_L' e^{-k\Lambda^2},
 \quad k=169/3200.
 \label{an03gesgv1:interiorspeedbound}
\end{equation}
The constants depend only on initial loss, Gaussian moments
and the ratio box, not on residual guards.
\end{lemma}
\begin{proof}
The two gate probabilities are at most
$e^{-\Lambda^2/2}$ and $e^{-\rho^2\Lambda^2/2}$.
For either cluster and any $i,j$, loss monotonicity and
Cauchy--Schwarz followed by Holder give
\[
 |\E_{P_k}[e_iX_j\one_A]|
 \le\|e_i\|_{P_k}(\E_{P_k}X_j^4)^{1/4}
                         \mathbb P_{P_k}(A)^{1/4}
 \le C_L e^{-k\Lambda^2}.
\]
All fourth moments are uniformly bounded for $q\le1$.
This proves the matrix estimate and then the three equations
in \eqref{an03gesgv1:exact} prove the speed estimate.
It strengthens the proposed $O(M_{\rm res}q^2)$ plus tails:
there is no need to discard either closed gate. In particular
no pre-$t_0$ cross-residual guard is needed.
\end{proof}

\begin{proposition}[Two initial strips isolate all Q3 recruitment]
\label{an03gesgv1:Q3split}
Let
\[
 \Lambda_q=\sqrt{3\ell/k},\qquad w_q=\arcsin(q\Lambda_q),\qquad
 \mathcal Q_{3,\mathrm{int}}^0
   =\{b\in Q3:\pi+2w_q\le b\le3\pi/2-2w_q\}.
\]
For sufficiently small $q$, on the entire interval $[0,t_*]$
these fixed initial labels stay in
$\pi+w_q<\theta<3\pi/2-w_q$, and
\begin{equation}
 |\theta(t,b)-b|+|\psi(t,b)-b|\le Cq^3\ell,\qquad
 m(t,b)=q^{10}e^{O(q^3\ell)}.
 \label{an03gesgv1:frozen}
\end{equation}
Consequently
\begin{equation}
 \int_{t_0}^{t_*}\frac{\mu_{\mathcal Q_{3,\mathrm{int}}^0}}q\dd t
 \le Cq^9\ell=o(1).
 \label{an03gesgv1:interiorcoststrong}
\end{equation}
More generally fix the subset of Q3 ancestry whose input
angle at $t_0$ lies in the same $2w_q$ interior. Its mass
satisfies the requested bound
\begin{equation}
 \int_{t_0}^{t_*}\frac{\mu_{\mathcal Q_{3,\mathrm{int}}(t_0)}}q\dd t
 \le C e^{2B}q\ell^3=o(1).
 \label{an03gesgv1:interiorcost}
\end{equation}
The set in this last display is fixed by its state at $t_0$;
it is not a moving angular mask.
\end{proposition}
\begin{proof}
MOMC's half-mass time and bounded $A_q$ give
$t_*\le C_{\rm dur}\ell$. Until the first exit from the
$w_q$ interior both projected margins exceed $q\Lambda_q$.
Lemma~\ref{an03gesgv1:interiorspeed} then bounds all speeds
by $Cq^3$, since $e^{-k\Lambda_q^2}=q^3$ exactly.
The available travel is $Cq^3C_{\rm dur}\ell=o(w_q)$,
because $w_q\ge q\Lambda_q$. It cannot consume the initial
angular buffer $w_q$. The first exit is therefore excluded.
Integrating the two angular speeds and the absolute radial
rate proves \eqref{an03gesgv1:frozen}, using the prescribed
initial mass $q^{10}$. Integration over the fixed label set
and over time proves \eqref{an03gesgv1:interiorcoststrong}.

Starting the same buffer argument at $t_0$ proves the analogous
retention for the entry-defined subset, independently of its
output angle. MOMC supplies its entry mass bound by the whole
left-ancestry mass $Cq^2\ell^2$. Its mass then changes by
$e^{O(q^3\ell)}$, giving $Cq\ell^3$ for the integrated
mass divided by $q$, and hence the weaker stated
$Ce^{2B}q\ell^3$ bound. Alternatively use
Lemma~\ref{an03gesgv1:lowernogain} on that retained set.
MOMC's mass bound alone would not locate these angles;
the preceding pre-$t_0$ speed proof supplies location for
the original-label interior.
\end{proof}

The only initial Q3 labels not covered by the stronger bound
are the two strips
\begin{equation}
 \mathcal S_\pi=[\pi,\pi+2w_q),\qquad
 \mathcal S_{3\pi/2}=(3\pi/2-2w_q,3\pi/2].
 \label{an03gesgv1:strips}
\end{equation}
Their widths are $2w_q=O(q\sqrt\ell)$. Neither strip has been
proved to join a specified family by $t_0$; naming them weak
or strong is a proposed classification, not a consequence of
the speed estimate. Endpoints are immaterial to the label
measure but can be included in their respective strip.

\section{Strip contracts in the form required by their consumers}
\label{an03gesgv1:stripcontracts}
The following are outstanding obligations, not extra hypotheses
of a claimed theorem. The two strips have not been folded into
the existing fixed partition.

\paragraph{The $\pi$ side: weak-family extension.}
An extension must first assign its labels to primary weak,
recruited secondary, or a separately paid remainder. For the
part assigned to $\Gp$, LCRC consumes at $t_0$
\[
 z>0,\quad |a|/z\le1/4,\quad qj/z\le C_0,\quad
 J(t_0)\le Cq^3,\quad E(t_0)=o(1),
\]
with the existing clock kept as a subset; enlarging the
clock is not required. For any part assigned to the recruited
group, the SEP/BULK consumer requires the raw radial submeasure
$\dd\eta_\pi=m(t_0,b)\dd\lambda_0$ and the actual-reference
displacement $\xi$ to satisfy
\begin{equation}
 \xi\le Z_\pi,\quad \eta_\pi(\mathbb R_+)\le Cq^3,\quad
 \eta_\pi\{\xi>u\}\le Cq^3\min\{1,(h_A/u)^{p_2}\},
 \quad h_A=A_q^dq^{10d-2s},\quad p_2=(1-3\lambda/2)/d.
 \label{an03gesgv1:pistripcontract}
\end{equation}
Here $A_q$ is bounded above and below, $d=(1-\lambda P_\rho)/2$
and $s=2d/(1-\lambda)$. Fixed constants suffice; BULK needs
the upper tail, not a two-sided asymptotic law. SEP presently
proves this contract only for its Q2 secondary group near
$\pi/2$, not for all labels near $\pi$. In particular a
primary weak label need not have bounded strong-chart
displacement. The proposed extension cannot assume that
every label in $\mathcal S_\pi$ belongs to that secondary group.
Any unassigned set needs an integrated outside-mass budget,
or retained negative-input geometry to use
\eqref{an03gesgv1:nogain}; strip membership alone is insufficient.

\paragraph{The $3\pi/2$ side: strong late-joiner extension.}
For the part assigned to the strong entry family, the useful
MOMC/PROFILE contract is a bounded actual-reference entry
chart and the mass-weighted upper tail
\begin{equation}
 \eta_{\rm late}\{\xi>u\}\le C
       \min\{1,(\epsilon_h/u)^{3/s}\},\qquad
 \epsilon_h\asymp A_q^d q^{10d-s},
 \quad \eta_{\rm late}=m(t_0,b)\dd\lambda_0,
 \label{an03gesgv1:latestripcontract}
\end{equation}
with a fixed total-mass bound and the same normalization as
the strong entry measure (MOMC's strong mass at $t_0$ is
$1/2$). If late labels are admitted to the safe core, their
entry displacement must also obey its cutoff. Otherwise
TAILREP applies to their complementary tail once the
all-endpoint gain is available. An equivalent alternative
is to prove directly the weighted tail moments
\[
 \int_{\{\xi>d_c\}}\max\{1,A\xi^2\}^{1+\nu}\dd\eta_{\rm late}
 \le C\epsilon_h^{3/s}d_c^{-3/s}
                  \max\{1,A d_c^2\}^{1+\nu}
\]
at the values of $A$ used by TAILREP. A label joining only
after $t_0$ is not covered by a fictitious charted entry at
$t_0$; it needs a rigorously matched arrival-time version or
a separate outside budget. We prove neither strip extension
in this increment.

\section{Outcome and explicit stopping point}
Task 1's matrix-value and splice calculations are closed,
but its fixed-constant conclusion stops at the visible
$\mathfrak X$ integral. Task 2's signed identities are closed,
including the removal of the unnecessary pointwise cap on
$N_-$. Its integrated-budget inference stops at
\eqref{an03gesgv1:missing}, or at
\eqref{an03gesgv1:missingnegative} if the full forcing budget
is explicitly guarded. Task 3 closes the constant-dependency
audit: fixed $M_{\rm res}$ dependence of the eventual envelope
constant is harmless in the specified consumers. Task 4's
lower-side, interior speed, pre-entry location and interior
mass bounds are closed; its two strip entry contracts remain
exactly as listed in Section~\ref{an03gesgv1:stripcontracts}.

The constants $C_R,C_L,C_L',k,R_c,C_{\rm dur}$ use only
initial loss, the ratio box and existing time/reference bounds.
$a_*,W_*,\varepsilon_q$ have the displayed residual-guard
dependence. $p_*,C_T,C_{\rm loc}$ depend on the fixed chart
and remainder budget. $B$ enters the envelope and no-gain
estimates exponentially. $K_Y$ enters only the diagnostic
cohort charge and the already specified downstream contracts.
No $q$-uniform value is assigned to $\int N_-$ or
$\mathfrak X$. The Gaussian choice of $\Lambda_q$ is an
explicit analytic formula, not a numerical tail test.

No joint weak-passage bootstrap, CLOCK passage theorem,
selected-field closure, selection, or consolidation is performed.
The completed gain integration and tail-moment results are
cited rather than reproved. Prior increments remain unchanged.

\begin{thebibliography}{9}\small
\bibitem{ROAD} Unified analytic theory roadmap v1.10:
S2 update after DECG and ledger items 15--17.
\bibitem{DECG} DECG, \texttt{AN03\_distance\_envelope\_and\_capped\_gain\_v1}:
Proposition 2.1, Lemma 3.1, Proposition 3.2, Corollary 3.3.
\bibitem{CLOCK} CLOCK, \texttt{AN03\_exact\_amplitude\_clock\_and\_c\_passage\_v1}:
c-form contract, forcing budget and exact clock.
\bibitem{CFR} CFR, \texttt{AN03\_c\_form\_field\_reduction\_v1}:
selected moments, non-core forcing, and negative-part lemma.
\bibitem{LCRC} LCRC, \texttt{AN03\_ledger\_cap\_residual\_closure\_v1}:
retained contracts and constant order.
\bibitem{TAIL} TAILREP and BULK: gain integration, tail moment,
secondary-energy implication and stopped clock.
\bibitem{ENTRY} MOMC and SEP: original strong half-mass entry,
left-ancestry mass and the recruited entry profile.
\bibitem{PROFILE} PROFILE: strong entry tail and exact
mass-charged outside field estimate.
\end{thebibliography}
\end{document}
